% \begin{savequote}
% \includegraphics[width=3.5cm]{./images/emhofer-josef-aass/taschenatlas_rettungsdienst.jpg}
% \end{savequote}
\def\ChapterCode{LER}
\chapter
[\CO{[LER]} Lernen]
{Lernen} 

\ChapterIntro
{%
Jeder Mensch hat im Laufe seines Lebens eigene Erfahrungen
mit Lernen gemacht und unter Umständen sogar seinen
optimalen Lernstil entdeckt. Aufgrund der Individualität der
Leser und Kursteilnehmer kann dieses Kapitel nicht das
\Speech{\enquote*{Rezept für Lernerfolg}} sein. Das soll es
auch nicht! In diesem Kapitel sollen lediglich Tipps und
weiterführende Ressourcen angeführt werden.
}
{%
\begin{description}
  \item[Maintainer:] Josef Emhofer
  \item[Autoren:] Diverse
  \item[Reviewer:] Standard-Reviewprozess
  \item[Version:] {\VarDocVersionTypeName} ({\VarDocVersionTypeDescription})
  \item[SHA1:] {\DocGitCommit}
\end{description}

{\TxTeilkapitelAASS}
}

\section{\HeadingKeyword{Strukturierung} des Lerninhalts}
Wenn dieses Buch als Lernunterlage für einen Kurs verwendet
wird, so können die Tipps dieses Abschnitts unter Umständen
hilfreich sein den umfangreichen Lehrstoff zu strukturieren.
Die Erfahrung zeigt, dass man leichter lernen kann, wenn man
die Details in einem großen Ganzen einordnen kann. Das
Gehirn kann dadurch Querverbindungen leichter herstellen.
Das vielzitierte \Speech{\enquote*{vernetzte Denken}}
gelingt dann effizienter. Es zahlt sich also aus sich einen
\EE{Überblick} zu verschaffen, bevor man sich auf ein Detail
stürzt.

\TextSlide{Eingrenzung des Lehrstoffs}{
Wenn man für eine bestimmte Prüfung lernt, ist es oft nicht
notwendig das gesamte Lehrwerk zu studieren.
Zwischenprüfungen werden meist nur über ein begrenztes
Stoffgebiet stattfinden. Bevor man zu lernen beginnt, sollte
man den relevanten Prüfungsstoff abgrenzen. Dazu gibt es
verschiedene Möglichkeiten:
\begin{itemize}
  \item Der Prüfungsstoff wird von den Lehrenden bekannt
  gegeben.
  \item Der Prüfungsstoff ergibt sich aus dem Curriculum
  (\Speech{\enquote*{Lehrplan}}) für die jeweilige
  Ausbildung.
  (Die Lehrpläne für Rettungs- und Notfallsanitäterkurse
  sind in der Ausbildungsverordnung zum Sanitätergesetz
  geregelt.)
  \item Der Prüfungsstoff wird durch die ausgeteilte
  Lehrunterlage definiert.
\end{itemize}
Im Zweifelsfall empfiehlt es sich (rechtzeitig) beim
Kursleiter zu informieren. 
}{
\begin{itemize}
  \item Bekanntgabe durch die Lehrenden
  \item Curriculum (\Speech{\enquote*{Lehrplan}})
  \item Lehrunterlage
\end{itemize}
}{}{}{}

\TextSlide{Strukturierung}{Zur Strukturierung stehen
verschiedene Möglichkeiten zur Verfügung:
  Die schnellste und einfachste Methode sich einen
  Überblick und eine erste Orientierung zu verschaffen ist
  sich ein paar Augenblicke mit dem 
  \EE{Inhaltsverzeichnis} zu beschäftigen. Durch die
  Gliederung des Inhalts eines (Lehr-) Buches kann man
  leicht eine (von vielleicht mehreren möglichen) Strukturen des
  Lehrstoffs gewinnen. Idealerweise versucht man die
  Attribute für die in diesem Buch gewählte Strukturierung
  zu erkennen.

Dieses Buch gliedert sich in folgende Ebenen:
\begin{itemize}
  \item \EE{Teile} sind Blöcke in welchen Kapitel
  zusammengefasst werden. So fasst z.\,B. der Teil
  \E{Medizinisch-Wissenschaftlicher Basisteil} jene Kapitel
  zusammen die hauptsächlich medizinisches Grundfachwissen
  vermitteln sollen. Der Zweck des Teils \E{Grundlagen} ist
  selbsterklärend. 
  \item \EE{Kapitel} sind große Einheiten eines spezifischen
  Themas. So gibt es im medizinischen Teil z.\,B. die
  Kapitel \nameref{chp:ANA} in welchem der Aufbau des
  menschlichen Körpers (ohne Krankheitsbilder) dargestellt
  ist oder das Kapitel \nameref{chp:KIN}
  welches speziell auf Kindernotfälle eingeht. Kapitel sind
  in diesem Buch zusätzlich zur Gliederungsnummer und zum
  Titel auch mit einem Code aus drei Buchstaben versehen.
  Dieser Code wird z.\,B. in Übersichtsgrafiken
  (\E{Leitsystem}) verwendet.
  \item Zwecks Übersicht und Strukturierung sind Kapitel
  selbst noch in \EE{Abschnitte} und \EE{Unterabschnitte}
  gegliedert.
\end{itemize}
Beispiel: Der \E{\nameref{topic:hitzekollaps}} ist der
Unterunterabschnitt \ref{topic:hitzekollaps} von
\begin{itemize}
  \item Kapitel \ref{chp:THE} \nameref{chp:THE}
  \begin{itemize}
    \item Abschnitt \ref{topic:Hitzeerkrankungen}
    \nameref{topic:Hitzeerkrankungen}
    \begin{itemize}
      \item Unterabschnitt \ref{topic:system-Hitzeschaeden}
      \nameref{topic:system-Hitzeschaeden}
    \end{itemize}
  \end{itemize}
\end{itemize}

Neben dem Inhaltsverzeichnis gibt auch der \EE{Stundenplan}
bzw. u.\,U. die Disposition des Kurses oder Seminars
Auskunft über die Gewichtung der einzelnen Teilinhalte und
somit über die Grobstruktur des Lehrstoffes.

In diesem Buch gibt es in einigen Kapiteln \EE{Querverweise}
und \EE{Leitsystem} die helfen sollen, das
aktuelle Kapitel in einen größeren Kontext einzubetten.

Eine sehr weit verbreitete Methode zum Strukturieren von
Inhalten ist die \EE{Mind-map}. Hier werden Zusammenhänge
von (Haupt-)Thema ausgehend durch Äste dargestellt. Je
weiter man ins Detail vor dringt, desto feiner wird die
Verästelung. Nach welchen Gesichtspunkten die Strukturierung
bzw. Gliederung durchgeführt wird bleibt dem Lernenden
selbst überlassen. Eine Möglichkeit ist natürlich die
Strukturierung/Gliederung des Buches zu übernehmen.
Abbildung \ref{fig:mindmap} ist die Strukturierung des
aktuellen Kapitels. Die Grafik wurde mit der freien Software
FreeMind erstellt. Wie man sieht können Themen/Begriffe
unterschiedlich formatiert und zwecks leichterer Merkbarkeit
mit Symbolen (Icons) versehen werden.


\Picture[H]
{Mind-map des aktuellen Kapitels}
{}
{\label{fig:mindmap}}
{}
{}
{./images/emhofer-josef-aass/Lernen.pdf}{}

}{
\begin{itemize}
  \item Inhaltsverzeichnis
  \item Stundenplan
  \item Querverweise
  \item Übersichtsgrafiken
  \item Mind-map
\end{itemize}
}{}{}{}

\section{\HeadingKeyword{Wissensquellen}}

\subsection{\HeadingKeyword{Fachbuchhandlungen}}

\Text{Facultas in der Berggasse}{%
Universitätsbuchhandlung für Medizin, 
Pharmazie, Gesundheitsberufe

\begin{compactitem}
  \item \EE{Anschrift}: Berggasse 2, A-1090 Wien
  \item \EE{Kontakt}: Tel +43\,1\,319\,53\,06; E-Mail
  \href{mailto:facultas@facultas.at}{\nolinkurl{facultas@facultas.at}}
  \item \EE{Öffnugszeiten}: \VonBis{Mo}{Fr} \VonBis{8:30}{18:30}
  Uhr, Sa \VonBis{9:00}{12:30} Uhr, Juli und August
  geänderte Öffnungszeiten
%   \item \EE{Anmerkungen}:
\end{compactitem}
}{}{}{}{}


\Text{Maudrich}{%
Universitätsbuchhandlung für Medizin, Pflege und Gesundheit

\begin{compactitem}
  \item \EE{Anschrift}: Spitalgasse 21a, A-1090 Wien
  \item \EE{Kontakt}: Tel +43\,1\,402\,47\,12;
  E-Mail: \href{mailto:medbook@maudrich.com}{\nolinkurl{medbook@maudrich.com}};
  WWW: \url{http://www.maudrich.com}
  \item \EE{Öffnugszeiten}: Öffnungszeiten: \VonBis{Mo}{Fr} \VonBis{8:30}{18:00} Uhr, Sa \VonBis{9:00}{12:00}
  Uhr, im Juli und August samstags geschlossen!
%   \item \EE{Anmerkungen}:
\end{compactitem}
}{}{}{}{}







\subsection{\HeadingKeyword{Bibliotheken}}

\Text{Universitätsbibliothek der Medizinischen Universität Wien}
{%
\begin{compactitem}
  \item \EE{Anschrift}: Allgemeines Krankenhaus (Neubau)
  Währinger Gürtel \VonBis{18}{20}, A-1097 Wien
  \item \EE{Kontakt}:
  \begin{inparadesc}
    \item[Tel.] +43\,1\,40160-26026;
    \item[WWW] \url{http://ub.meduniwien.ac.at};
    \item[E-Mail für Fernleihebestellungen] ubmed-fernleihe@meduniwien.ac.at
  \end{inparadesc}
  \item \EE{Öffnugszeiten}:
  \begin{inparadesc}
    \item[\VonBis{Mo}{Fr}] \VonBis{8:00}{20:00},
    \item[Sa] \VonBis{9:00}{17:00}
    \item[Sonst.:] In den Ferien gelten geänderte
    Öffnungszeiten. Die Schalter (Kassa, Entlehnung,
    Ausweise, \ldots) haben abweichende Öffnungszeiten.
  \end{inparadesc}
  \item \EE{Anmerkungen}: Voraussetzung für die
  Entlehnberechtigung ist eine gültige Bibliothekskarte der
  Universitätsbibliothek der Medizinischen Universität Wien.
  Dies gilt auch für Personen, die schon an anderen
  Universitätsbibliotheken entlehnberechtigt sind.
\end{compactitem}
}{}{}{}{}















% Hauptbibliothek
% 
% Nationalbobliothek


\subsection{\HeadingKeyword{Internet}}

% Youtube



% Fischer Lexikon Medizin

% EKG-Kurs

\begin{itemize}
  \item % Online-EKG-Kurs: \url{http://www.ekg-online.de}
  \fullcite{www.ekg-online.de}
%   \item Birth of Baby (Vaginal Childbirth)
%   \url{http://www.youtube.com/watch?v=Xath6kOf0NE}
  \item \fullcite{youtube:VaginalChildbirth}
\end{itemize}


% Wikipedia

% DocCeck-Lexikon?

\clearpage
\section{\HeadingKeyword{Literatur}}

\TextSlide{Bücher -- Empfehlung der Redaktion}{
\begin{itemize}
  \item \EE{Lehrbuch für
  Fortgeschrittene}:
  {\fullcite{RedelsteinerHRNS2}}
  \item Lehrbuch und Atlas für die Berufe im
  Gesundheitswesen: 
  
  \fullcite{MenschKoerperKrankheit6}
  \item \EE{Fachbuch} für speziell an Medizin
  interessierte, mit guter, kurzer Einführung in anatomische
  und physiologische Grundlagen am Anfang eines jeden Kapitels:
  
  \fullcite{BLIM4}
  \item \EE{Taschenbuch}, jackentaschentauglich:
  
  \fullcite{BoehmerTaschRett6}
  \item \EE{Umfangreiches
  Anästhesie-Fachbuch} in 2 Bänden. Der Striebel ist ein
  herausragendes Buch: Didaktisch hervorragend aufbereitet,
  reich bebildert und illustriert und konkrete, verständliche
  Formulierungen.
  
  Das Werk richtet sich zwar vorrangig an Ärzte, für fortgeschrittenes
  nichtärztliches Fachpersonal stellt es jedoch eine ebenso
  wertvolle Lern- und Wissensquelle dar:
  
  \fullcite{StriebelAnaesthesie2}
\end{itemize}
}{
\begin{itemize}
  \item Redelsteiner: \emph{Das Handbuch für Notfall- und Rettungssanitäter}
  \item Huch: \emph{Mensch Körper Krankheit}
  \item Renz Polster: \emph{Basislehrbuch Innere Medizin}
  \item Böhmer: \emph{Taschenatlas Rettungsdienst}
\end{itemize}
}{}{}{}

\TextSlide{Bücher -- weitere Empfehlungen}{%
\begin{itemize}
  \item \EE{Kurzgefasste Basics} der Notfallmedizin auf 141
  Seiten:
  
  \fullcite{Helfen2008}
\end{itemize}
}{
\begin{itemize}
  \item Helfen: \emph{BASICS
Notfall- und Rettungsmedizin}
\end{itemize}
}{}{}{}


\TextSlide{Zeitschriften -- Empfehlung der Redaktion}{% 

\begin{compactdesc}
  \item[Rettungsdienst] S+K-Verlag. 
  
  Erscheint einmal
monatlich und durchgehend vierfarbig mit Magazinteil,
beinhaltet aktuelle Informationen aus dem Bereich der
präklinischen Notfallmedizin, Fallberichte, zertifizierte
Fortbildung, Notfallpraxis, Recht, Berufspolitik u.\,a.

Homepage: \url{http://www.skverlag.de/zeitschriften/rettungsdienst.html}
\end{compactdesc}
}{
\begin{itemize}
  \item \emph{Rettungsdienst}, S+K-Verlag
\end{itemize}

}{}{}{}


% \section{Prüfungen}
% 
% \subsection{Anweisungen}
% 
% \begin{enumerate}
%   \item Am Begin eines Fallbeispieles erhalten Sie die Informationen zur
%   Situation und Szene.
%   \item Dokumentieren Sei jeden Befund (laut ansagen)! Erst dann passt der
%   Prüfer den Befund an das Prüfungsbeispiel an.
%   \item Die simulierte Funkkennung lautet \TakeNotesMedium . 
%   \item Telefongespräche, z.B. mit der Leitstelle, sind mit dem simulierten
%   Telefon zu führen. Die gewählte Rufnummer ist laut anzusagen.
% \end{enumerate}
% 
% \subsection{Globalskalen}
% 
% \begin{WidePar}
% 
% \begin{tabulary}{\linewidth}{LLLLL}
% 1 & 2 & 3 & 4 & 5
% \\\addlinespace\hline\addlinespace
% \multicolumn{5}{c}{\TabHeading{Untersuchungsablauf}}
% \\\addlinespace
% maximale Effizienz im Arbeitsablauf,klare Bewegungsökonomie (müheloser Fluss
% von einer Bewegung zur nächsten) 
% 
% Durchführung weitgehend vollständig\,/\,in weitgehend korrekter Reihenfolge
% &
% effizienter Arbeitsablauf überwiegt, vereinzelt unnötige Bewegungen
% 
% Durchführung weitgehend vollständig\,/\,in weitgehend korrekter Reihenfolge
% durchgeführt, auf Widrigkeiten im Ablauf wird korrekt reagiert
% &
% teilweise effizienter Arbeitsablauf, vermehrt unnötige Bewegungen
% 
% Durchführung ausreichend vollständig\,/\,in ausreichend korrekter Reihenfolge
% durchgeführt; auf Widrigkeiten im Ablauf wird nur teilweise korrekt reagiert
% &
% Ineffizienter Arbeitsablauf, viele unnötige Bewegungen und/oder Unsicherheiten
% bezüglich der nächsten Bewegung
% 
% Durchführung in relevanten Teilbereichen unvollständig oder in falscher
% Reihenfolge; auf Widrigkeiten im Ablauf kann nicht oder nur unangemessen
% reagiert werden
% &
% Durchführung patienten-, eigen- oder drittgefährdend oder für das
% Patienten-Outcome wesentliche Untersuchungen wurden nicht durchgeführt
% \\\addlinespace\hline\addlinespace
% \multicolumn{5}{c}{\TabHeading{Instrumentenhandhabung}}
% \\\addlinespace
% Flüssige, routinierte Bewegungen mit den Instrumenten
% 
% gschickte Instrumentenanwendung
% &
% kompetente Verwendung der Instrumente überwiegt
% 
% vereinzelt ungeschickte Instrumentenanwendung
% &
% unroutinierte Verwendung der Instrumente
% 
% vermehrt ungeschickte Instrumtenenanwendung
% &
% unangemessene Verwendung der Instrumente
% 
% deutlich ungeschickte Instrumentenanwendung
% &
% patienten-, eigen- oder drittgefährdend
% \\\addlinespace\hline\addlinespace
% \multicolumn{5}{c}{\TabHeading{Behandlung und Maßnahmen}}
% \\
% Optimale Behandlung
% 
% Durchführung weitgehend vollständig\,/\,in weitgehend korrekter Reihenfolge
% &
% kompetente Behandlung überwiegt, vereinzelt unnötige Bewegungen
% 
% Behandlung weitgehend vollständig\,/\,in weitgehend korrekter Reihenfolge
% durchgeführt, auf Widrigkeiten wird korrekt reagiert
% &
% teilweise kompetente Behandlung, teilweise unnötige Maßnahmen
% 
% Behandlung ausreichend vollständig\,/\,in ausreichend korrekter Reihenfolge
% durchgeführt; auf Widrigkeiten wird nur teilweise korrekt reagiert
% &
% Ineffiziente Behandlung oder falsche Behandlung \E{ohne Einfluss auf das
% Patienten-Outcome}, viele unnötige Maßnahmen und/oder Unsicherheiten bezüglich
% der nächsten Maßnahme
% 
% Durchführung in relevanten Teilbereichen unvollständig oder in falscher
% Reihenfolge; auf Widrigkeiten kann nicht oder nur unangemessen
% reagiert werden
% &
% Durchführung patienten-, eigen- oder drittgefährdend oder für das Patienten-Outcome
% wesentliche  Maßnahmen
% wurden nicht durchgeführt 
% \\\addlinespace\hline\addlinespace
% \multicolumn{5}{c}{\TabHeading{Haltung gegenüber dem Patienten}}
% \\\addlinespace
% aufmerksam, warm, entspannt, umgänglich, geht auf die Bedürfnisse des Pateinten
% ein
% 
% optimales Vorstellen, optimale Erklärung der Prozedur, optimales Verabschieden
% &
% Angemessene, aber eher neutrale Grundhaltung
% 
% Korrektes Vorstellen, Korrekte Erklärung der Prozedur, Korrektes Verabschieden
% &
% Unbeholfen oder unsicher oder zögerlich oder überheblich oder kühl oder
% gehetzt oder abrupt
% 
% Mängel beim Vorstellen und/oder bei der Erklärung der Prozedur und/oder beim
% Verabschieden
% &
% unpersönlich, geht nicht auf die Bedürfnisse des Patienten ein
% 
% kein oder unzureichendes Vorstellen, keine oder unzureichende Erklärung der
% Prozedur, kein oder unzureuchendes Verabschieden
% &
% beleidigend oder prvozierendes Auftreten oder patienten-, eigen- oder
% drittgefährdend 
% \\\addlinespace\hline\addlinespace
% \multicolumn{5}{c}{\TabHeading{Gesprächsverlauf}}
% \\
% maximale Effizienz im Einholen der Informationen
% 
% Idealer Wechsel zwischen offenem und geschlossenem Frgatentyp
% 
% strukturierter, aber flexibler Gesprächsfluss
% &
% Überwiegend Effizienz iom Einholen der Informationen
% 
% angemessener Wechsel zwischen offenem und geschlossenem Fragentyp
% 
% überwiegend strukturierter, aber flexibler Gesprächsfluss
% &
% teilweise Effizenz im Einholen der Informationen
% 
% kein Wechsel des Fragentyps und/oder unangemesenes Verwenden
% eines Fragentyps
% 
% eher unstrukturierter und/oder starrer Gesrächsfluss
% &
% Ineffizienz im Einholen der Information (z.B. Unsicherheit über die als
% nächstes zu stellende Frage)
% 
% Unangemessener Fragentyp überwiegt (Suggestivfragen)
% 
% unstrukturierter und/oder starrer Gesprächsfluss
% &
% erheben von falschen Informationen z.B. durch Suggestivfragen
% 
% patienten-, eigen- oder drittgefährdend oder für das Patienten-Outcome
% wesentliche Informationen wurden nicht eingeholt
%  \\\addlinespace\hline
% \end{tabulary}
% 
% 
% \end{WidePar}

\nocite{Steinert2008}
\ChapterEnd